% =========================================
% ACRÓNIMOS Y GLOSARIO
% =========================================
% 
% Para usar este archivo:
% 1. Descomentar en TFG.tex: % =========================================
% ACRÓNIMOS Y GLOSARIO
% =========================================
% 
% Para usar este archivo:
% 1. Descomentar en TFG.tex: % =========================================
% ACRÓNIMOS Y GLOSARIO
% =========================================
% 
% Para usar este archivo:
% 1. Descomentar en TFG.tex: % =========================================
% ACRÓNIMOS Y GLOSARIO
% =========================================
% 
% Para usar este archivo:
% 1. Descomentar en TFG.tex: \input{etc/acro}
% 2. Descomentar en TFG.tex: \printglossary[type=\acronymtype,title={Abreviaciones y Acrónimos}]
%
% Comandos útiles:
% \gls{acronimo}      - Primera vez: forma completa, siguientes: forma corta
% \acrshort{acronimo} - Solo la forma corta (ej: LLM)
% \acrlong{acronimo}  - Solo la forma larga (ej: Large Language Model)
% \acrfull{acronimo}  - Forma completa (ej: Large Language Model (LLM))

\usepackage[toc,acronym]{glossaries}
\makeglossaries

% Define tus acrónimos aquí
% Ejemplo:
% \newacronym{llm}{LLM}{Large Language Model}
% \newacronym{rag}{RAG}{Retrieval-Augmented Generation}
% \newacronym{ia}{IA}{Inteligencia Artificial}
% \newacronym{pln}{PLN}{Procesamiento de Lenguaje Natural}
% \newacronym{sql}{SQL}{Structured Query Language}
% \newacronym{api}{API}{Application Programming Interface}











% 2. Descomentar en TFG.tex: \printglossary[type=\acronymtype,title={Abreviaciones y Acrónimos}]
%
% Comandos útiles:
% \gls{acronimo}      - Primera vez: forma completa, siguientes: forma corta
% \acrshort{acronimo} - Solo la forma corta (ej: LLM)
% \acrlong{acronimo}  - Solo la forma larga (ej: Large Language Model)
% \acrfull{acronimo}  - Forma completa (ej: Large Language Model (LLM))

\usepackage[toc,acronym]{glossaries}
\makeglossaries

% Define tus acrónimos aquí
% Ejemplo:
% \newacronym{llm}{LLM}{Large Language Model}
% \newacronym{rag}{RAG}{Retrieval-Augmented Generation}
% \newacronym{ia}{IA}{Inteligencia Artificial}
% \newacronym{pln}{PLN}{Procesamiento de Lenguaje Natural}
% \newacronym{sql}{SQL}{Structured Query Language}
% \newacronym{api}{API}{Application Programming Interface}











% 2. Descomentar en TFG.tex: \printglossary[type=\acronymtype,title={Abreviaciones y Acrónimos}]
%
% Comandos útiles:
% \gls{acronimo}      - Primera vez: forma completa, siguientes: forma corta
% \acrshort{acronimo} - Solo la forma corta (ej: LLM)
% \acrlong{acronimo}  - Solo la forma larga (ej: Large Language Model)
% \acrfull{acronimo}  - Forma completa (ej: Large Language Model (LLM))

\usepackage[toc,acronym]{glossaries}
\makeglossaries

% Define tus acrónimos aquí
% Ejemplo:
% \newacronym{llm}{LLM}{Large Language Model}
% \newacronym{rag}{RAG}{Retrieval-Augmented Generation}
% \newacronym{ia}{IA}{Inteligencia Artificial}
% \newacronym{pln}{PLN}{Procesamiento de Lenguaje Natural}
% \newacronym{sql}{SQL}{Structured Query Language}
% \newacronym{api}{API}{Application Programming Interface}











% 2. Descomentar en TFG.tex: \printglossary[type=\acronymtype,title={Abreviaciones y Acrónimos}]
%
% Comandos útiles:
% \gls{acronimo}      - Primera vez: forma completa, siguientes: forma corta
% \acrshort{acronimo} - Solo la forma corta (ej: LLM)
% \acrlong{acronimo}  - Solo la forma larga (ej: Large Language Model)
% \acrfull{acronimo}  - Forma completa (ej: Large Language Model (LLM))

\usepackage[toc,acronym]{glossaries}
\makeglossaries

% Define tus acrónimos aquí
% Ejemplo:
% \newacronym{llm}{LLM}{Large Language Model}
% \newacronym{rag}{RAG}{Retrieval-Augmented Generation}
% \newacronym{ia}{IA}{Inteligencia Artificial}
% \newacronym{pln}{PLN}{Procesamiento de Lenguaje Natural}
% \newacronym{sql}{SQL}{Structured Query Language}
% \newacronym{api}{API}{Application Programming Interface}










